\documentclass[10pt]{article}
\usepackage[margin=0.9in]{geometry}
\usepackage[T1]{fontenc}
\usepackage[utf8]{inputenc}
\usepackage{amsmath,amssymb,booktabs,enumitem}
\usepackage{graphicx}
\usepackage{xcolor}
\usepackage{hyperref}
\usepackage{natbib}
\setlength{\parindent}{0pt}
\setlength{\parskip}{0.5em}

\newcommand{\remark}[1]{{\color{blue}\textbf{Remark.} #1}}

% \title{The Emergence of Axiomatic Attention Patterns During LoRA Fine Tuning of Rerankers}
\title{The Emergence of Relevance Through Axiomatic Attention Patterns During LoRA Fine Tuning}
\author{Matthew Perlman, Atharva Nijasure, James Allan}
\date{April 2026}

\begin{document}

\maketitle

\begin{abstract}
LoRA and fine-tuning are standard tools for adapting large language models to reranking, but it remains unclear where task-specific relevance behavior is learned and what internal changes accompany that learning.
Through ablation and attention-mass analyses, this paper studies whether reranking gains are localized to a narrow subset of layers and whether those gains are associated with interpretable relevance-oriented attention patterns such as lexical matching, semantic matching, rarity sensitivity, and query-document interaction.
We find that the strongest learned relevance signals appear in a compact mid-network region and are strongly associated with feature-level attention changes that resemble classical IR intuitions, suggesting both a mechanistic account of learned relevance and the possibility of more targeted LoRA adaptation.
\end{abstract}

\section{Introduction}

LoRA has become a practical default for adapting large language models to domain-specific tasks, including neural reranking \citep{hu2022lora}.
For reranking in particular, this raises a concrete interpretability question: after fine-tuning, where in the network is task-specific relevance behavior actually implemented, and what internal patterns change with it?

This paper studies that question using two complementary analyses.
First, we perform layer, window, and keep-window ablations to localize where LoRA updates matter for ranking performance.
Second, we analyze attention mass over interpretable token-feature pairs to characterize what kinds of relevance-oriented behavior emerge in those same regions.

\remark{One possible way to keep the introduction focused is to make a single clear promise: this paper is not mainly about general LLM interpretability or head cataloging, but rather about whether reranking-relevant LoRA adaptation is localized, whether that localized region corresponds to particular attention-head behaviors, and whether those behaviors align with interpretable, relevance-oriented attention patterns.}

\section{Research Questions}

\remark{It may help to include this section explicitly, since it gives the paper a cleaner scope and keeps the introduction from sounding too broad.}

\begin{enumerate}[label=\textbf{RQ\arabic*.}, leftmargin=*, itemsep=4pt]
  \item Where in the transformer do LoRA updates contribute most to reranking performance?
  \item Are those performance-critical regions associated with specific relevance-oriented attention patterns, such as exact lexical matching, semantic matching, rarity sensitivity, or query-document interaction?
  \item Does the overlap between performance-critical regions and feature-level attention shifts support a mechanistic interpretation of learned relevance?
\end{enumerate}

\remark{If you decide to make the framing even tighter for an EMNLP-style paper, one option would be to keep only RQ1 and RQ2 in the main text and leave RQ3 implicit in the discussion.}

\section{Related Work}

This work sits at the intersection of neural reranking, parameter-efficient adaptation, and mechanistic analysis of relevance.
On the retrieval side, \citet{ma2024rankllama} establish the RankLLaMA setting that motivates this paper: a decoder-based large language model adapted for multi-stage text retrieval.
The present work does not introduce a new reranker architecture or training objective; instead, it asks what the existing LoRA adaptation is functionally changing inside such a reranker.

Among the closest prior work, \citet{nijasure2025probing} show that classical IR-relevant signals can be recovered from internal representations of ranking LLMs.
That paper is close in spirit because it also asks whether retrieval-relevant structure appears inside the model, but the emphasis is different.
Their focus is representational: whether features can be decoded from hidden states.
The focus here is more specifically on adaptation and localization: which layers appear to acquire task-specific relevance behavior during LoRA fine-tuning, and whether those layers are associated with interpretable attention shifts.

The relevance-mechanism perspective is also related to \citet{lu2025pathway}, who show that cross-encoders can implement a semantic analogue of BM25.
That paper is a useful reference point because it argues that classical IR-style relevance behavior can emerge inside neural rankers, but the present project differs in both architecture and goal.
Lu et al.\ study a cross-encoder and reconstruct a pathway for semantic BM25-like behavior, whereas this paper studies a decoder-based reranker and asks how LoRA fine-tuning changes the internal distribution of relevance-oriented behavior across layers.

More broadly, the idea that useful task adaptation may concentrate in a relatively small subset of parameters or components is consistent with sparsity-oriented work such as the lottery ticket hypothesis \citep{frankle2019lottery}.
This paper does not attempt to prove a lottery-ticket style claim.
Instead, it uses that broader intuition as motivation for asking whether reranking gains from LoRA are similarly localized, and whether that localization can be interpreted in terms of token-level relevance behavior.

\remark{It may be useful to keep this related-work section distinct from the synthesis-project draft. Rather than reusing the longer probing/intervention framing, this version could stay centered on LoRA adaptation, sparse task specialization, and IR relevance mechanisms.}

\section{Ablation Experiments}

The first experimental question is where LoRA fine-tuning matters for reranking.
To answer that, we compare the fully fine-tuned reranker against variants in which specific LoRA-affected layers or contiguous windows are restored to their base-model values.
This lets us estimate which parts of the network are most responsible for the observed ranking gains.

\subsection{Methodology}

\remark{It would probably help to define the intervention very precisely here: the experiment restores selected fine-tuned attention weights to their base-model values, rather than deleting layers or zeroing activations. That distinction matters for the causal interpretation.}

\remark{A concise way to organize this subsection might be to list the intervention families first: single-layer ablation, keep-single-layer, window ablation, and keep-window ablation. From there, you could explain that window and keep-window are the main evidence because they appear more stable than the single-head or single-layer results?}

\subsection{Results}


\section{Feature Attention Mass Experiments}

The second experimental question is what kinds of token-level interactions become more prominent in the regions identified by ablation.
To study that, we aggregate attention mass over interpretable token-pair features including lexical matching, semantic similarity, rarity, and query-document interaction structure.

\subsection{Methodology}



\subsection{Results}


\section{Discussion}



\section{Future Work}

\remark{Should we use this as a project tracker? and big picture future work we can write just before the paper submission? }

\remark{One possible ordering for future work is to start with what most directly strengthens the current paper: stronger evaluation and stability checks first, then injection or causal patching, and only after that the larger circuit-discovery ideas.}

\begin{enumerate}[label=\textbf{\arabic*.}, leftmargin=*]
  \item Evaluate the same interventions on stronger reranking settings with realistic candidate sets rather than lightweight random-negative evaluation.
  \item Test sufficiency more directly through injection-style experiments that add selected LoRA deltas back into a base model.
  \item Refine the feature analysis by isolating which projections, heads, or submodules are responsible for the observed relevance-oriented attention shifts.
\end{enumerate}

\remark{If space becomes tight, it may be enough for the future-work section to briefly mention causal patching and module-level injection without expanding into a separate agenda. Keeping it closely tied to the present paper's claims may make the ending feel more focused.}

\bibliographystyle{plainnat}
\bibliography{main_refs}

\end{document}
